\newpage
%\TOCadd{ExecutiveSummary}

\phantomsection

\section*{}

\begin{flushright}
	142 Flower Street\\
	Brisbane, Q 4013\\
	Tel.\ 0400 012 299\\
	\medskip
	\today
\end{flushright}
\begin{flushleft}
  Associate Professor Geoffrey Walker\\
  School of Science and Engineering\\
  Queensland University of Technology\\
  Brisbane City, Q 4000\\
  \bigskip\bigskip
  Dear Professor Walker,
\end{flushleft}

In accordance with the requirements of the degree of Bachelor of
Engineering / Bachelor of Finance (IX28) in the division of Computer Systems
Engineering / Electrical and Electronic Engineering, I present the
following thesis entitled Feasibility Study of Low Voltage Direct Current Power Distribution.  This work was performed under the supervision of Associate Professor Geoffrey Walker.

I declare that the work submitted in this thesis is my own, except as
acknowledged in the text and footnotes, and has not been previously
submitted for a degree at Queensland University of Technology or any other
institution.

\begin{flushright}
	Yours sincerely,\\
	\medskip
	\begin{figure}[H]
	\hfill\includegraphics[width = 40mm]{images/Signature.JPG}%\hspace*{\fill}
	\end{figure} 
	\medskip
	DAVID PETRIE.
\end{flushright}

\newpage

\section*{Acknowledgements}


\paragraph{}
\textbf{Associate Professor Geoffrey Walker} for your guidance, support and patience through the completion of this project. 

\paragraph{}
My fellow students \textbf{Ash Abdullrabzak} and \textbf{Niroj Gurung} for their valuable discussions throughout our three projects. 

\paragraph{}
\textbf{Disa}, for her constant love and support through not only this project but my entire degree.    

\paragraph{}
Past Queensland University of Technology engineering student \textbf{Steven Donohue} for his valuable prior thesis on 48 V direct current home power systems.  

\paragraph{}
The Electrical Engineering team in Buildings at Aurecon, Brisbane for their technical advice and assistance in the final half of my project.  

\paragraph{}
\textbf{Geoffrey Woods} and \textbf{Norman Higgins} from Queensland University of Technology's Building Management Services for providing technical drawings and electricity data.

\paragraph{}
\textbf{Abhishek Bhasin} and \textbf{Isaac Linett} from Q Electrical for their technical assistance in the first half of my project.



\newpage


\section*{Abstract}

\paragraph{}
This design project aims to explore the issue of incorporating low or extra low voltage direct current (DC) power distribution systems to multi-residential and commercial buildings. Low voltage DC power is prevalent in telecommunications systems but has the possibility of being used for a variety of other applications to reduce electricity costs and improve efficiency. Due to the large costs involved in building power systems, the majority of this design project will be completed through software simulations and hand calculations. If these methods cannot be utilised, either physical testing or research will be completed as substitute. If successful, a more efficient and cheaper alternative to running simple electronics from Alternating Current (AC) mains will be created. This will allow for a portion of building's electricity demands to be self sufficient and not rely on the grid for continual supply. 

\paragraph{}
To complete this task, the project was separated into individual questions that are interrelated and help to conclude whether the project will be successful or not. These questions will cover finer details of the project including the chosen DC voltage level, whether photovoltatics can be used, if lighting requirements will be met, structural design and safety mechanisms for voltage loss reduction and finally whether DC could reasonably replace AC for portions of commercial building power systems. Simulations were completed using System Advisor Model (SAM) for photovoltaics, Dialux4.12 for lighting, AutoCad for structural layouts and Microsoft Excel for all calculations and efficiency modelling.  

\paragraph{}
***FILL WITH RESULTS SUMMARY***
 
\newpage
